\section{Protocol-integration gap}
\label{sec:protocol-gap}

\subsection{The issue}

The encoder + \code{LinearCode<Zt>} + \code{ZipTypes} machinery is
complete and exercised by tests. But the existing Zip$^+$
prove/verify path makes an implicit assumption that does not hold
for our \code{CombR}.

Concretely, at \code{phase\_prove.rs:280} the prover computes
\[
  b_j \;=\; \sum_i \mathrm{MBSInnerProduct.inner\_product\_field}
            \bigl( \text{row}_{ij}, q_1, 0_F \bigr),
\]
where $\text{row}_{ij}$ is a slice of \code{Zt::CombR} elements
and $q_1$ is a slice of \code{F} (the verifier-side prime field)
elements. This requires
\[
  F \,:\, \mathrm{FromWithConfig}\!\left\langle\,
    \&\,\text{Zt::CombR}\,
  \right\rangle,
\]
i.e.\ a mapping
\[
  \mathrm{Zt::CombR} \;\longrightarrow\; F.
\]
For the existing IPRS / RAA codes the \code{CombR} is a scalar
\code{Int<N>} and the mapping is the standard ``integer mod $p$'' lift.
For our \code{CombR = DensePolynomial<Int<M>, 32>} no such mapping
exists out of the box, and \code{cargo build} explicitly reports:

\begin{quote}\small
\code{error[E0277]: the trait bound `BoxedMontyField: FromWithConfig<\&DensePolynomial<Int<8>, 32>>` is not satisfied}
\end{quote}

\subsection{The intended fix (clarified after the fact)}

The original design described in the plan handed the entire
prove/verify off to ``run unchanged''. After surfacing the gap and
asking the user, the intended design is that the prover's
combined-row vector $b$, the evaluation $\mathrm{eval}_F$, and the
tensor $(q_0, q_1)$ are \emph{not} prime-field-valued. Instead,
they live in
\[
  F[X]/\flift,
\]
the same polynomial ring shape as \code{CombR} but with prime-field
coefficients in place of $\Z$-coefficients. The map
$\Rlift \to F[X]/\flift$ is the coefficient-wise mod-$p$ reduction
$\Z \twoheadrightarrow \F_p$, which is a ring homomorphism, so it
preserves the protocol's algebraic identities.

\subsection{What it would take to land this}

Implementing the polynomial-valued prove/verify is a
protocol-level refactor, not a localized change. At minimum it
requires:

\begin{enumerate}
  \item \textbf{Wrap $F[X]/\flift$ as a new ``verifier-side scalar''
    type} that exposes whatever subset of \code{PrimeField} the
    prove/verify code actually uses. If $\flift \bmod p$ is
    irreducible, this is a genuine field extension $\F_{p^{32}}$;
    if not, it is a ring (and the protocol must avoid relying on
    inversion).
  \item \textbf{Generalize the prove/verify generic} from
    \code{F: PrimeField} to a weaker trait satisfied by both
    \code{BoxedMontyField} (the existing IPRS path) and the new
    $F[X]/\flift$ wrapper.
  \item \textbf{Update the transcript layout} so that the
    combined-row vector $b$ serializes as polynomial-valued entries
    rather than \code{write\_field\_elements}.
  \item \textbf{Decide on verifier-challenge shape} ---
    most cleanly the row-combination coefficients and tensor
    coordinates are also drawn from $F[X]/\flift$.
  \item \textbf{Soundness analysis.} The Zip$^+$ proximity
    argument leans on $F$'s prime-field structure; an argument
    over $F[X]/\flift$ needs a separate pass, plausibly via
    the same Schwartz--Zippel-style argument applied to
    $\F_{p^{32}}$ if $\flift$ is irreducible mod $p$.
\end{enumerate}

This is left as follow-up work. The encoder + ZipTypes side
delivered here is reusable verbatim once the polynomial-valued
prove/verify lands. Until then, the bench harness exercises only
\code{encode\_rows} / \code{commit} / Merkle, which do not require
the field projection.
